\section{Conclusion}\label{sec:conclusion}
ELEC-VQA combines 10,201 engineering QA records, 1,019 schematic configurations, compact evidence and source-disjoint evaluation. The comparative study, including provisional supplementary adjudication, demonstrates substantial domain-adaptation gains, but distinguishes answer accuracy from supported reasoning. Answer-only adaptation achieves the strongest concise-answer result; evidence-guided adaptation yields higher joint correctness and more reliable support conditional on correct answers. Gains concentrate on measurement planning and selected multi-condition tasks. Controlled-change pairs expose complementary failures in irrelevant-state invariance and visual topology sensitivity. These results motivate evidence supervision that models both relevant and irrelevant electrical changes, rather than assuming that longer explanations universally improve understanding. Broader drawing coverage, independent reference validation and multiple training seeds are needed to establish generality.
